\section{Solutions for Chapter~\ref{ch:finite_closure}}

\paragraph{Exercise~\ref{ex:c-closure-1}: all negative modes.}
\phantomsection
\label{sol:ex:c-closure-1}
Translation covariance gives $(Da)_{(n)}=-n a_{(n-1)}$.
For example $(D^2a)_{(-1)}=(Da)_{(-2)}=2a_{(-3)}$.
More generally each differentiation supplies the next factor
$1,2,\ldots,r-1$, giving
$(D^{r-1}a)_{(-1)}=(r-1)!a_{(-r)}$. Thus negative modes are
normal products with translated left inputs. Once translation, normal
products, and nonnegative modes preserve a space, every mode preserves
it, which is exactly vertex-subalgebra closure.



\paragraph{Exercise~\ref{ex:c-closure-2}: the missing lowest line.}
\phantomsection
\label{sol:ex:c-closure-2}
Choose weight bases $v_0,\ldots,v_{m-1}$ and
$w_0,\ldots,w_{n-1}$ ordered from lowest to highest. The raising
image in the first factor contains every $v_i$ with $i>0$, and
the second contains every $w_j$ with $j>0$. Their sum therefore
contains all tensor basis vectors except $v_0\otimes w_0$.
That vector cannot be in the image because neither raising operator
can reach a joint-lowest weight. The quotient is one dimensional.
Equivalently, normalize the basis so that $E_Lv_i=v_{i+1}$
below the top and $E_Rw_j=w_{j+1}$. Every basis tensor except
$v_0\otimes w_0$ then has an explicit preimage under one of the
two operators. This is a quotient of vector spaces; it need not
be a representation of the lowering operators.



\paragraph{Exercise~\ref{ex:c-closure-3}: counting the raising quotient.}
\phantomsection
\label{sol:ex:c-closure-3}
The summands whose left representation has dimension one contribute
\[
 3+1+5+2+3+1+2+1=18
\]
irreducible constituents and
\[
 9+5+35+18+33+13+30+19=162
\]
dimensions. Those whose left representation has dimension two
contribute $1+1+2+2+2+1+1=10$ constituents and
$4+8+24+32+40+24+28=160$ dimensions. Thus $P$ has dimension
$322$ and $28$ constituents, with multiplicity. Each constituent
contributes precisely one joint-lowest vector to the quotient by
$E_LP+E_RP$, by Exercise~\ref{ex:c-closure-2}. The image therefore
has dimension $322-28=294$. Multiplicity copies require separate
seeds even when their representation types agree.



\paragraph{Exercise~\ref{ex:c-closure-4}: specialization can lose rank.}
\phantomsection
\label{sol:ex:c-closure-4}
The matrix $\left(\begin{smallmatrix}1&0\\0&p\end{smallmatrix}\right)$
has rank two over $\Q$ and rank one over $\mathbb F_p$.
A nonzero minor modulo $p$ was nonzero over $\Q$, but a vanishing
minor modulo $p$ may merely have been divisible by $p$. The scalar
residual $p$ is the simplest example of a nonzero integer residual
that reduces to zero. Thus lower rank certificates and exact
zero-residual certificates require different evidence.
For rational matrices, first clear denominators and choose a prime
which does not divide them. Without that condition, the proposed
specialization might not even be defined.



\paragraph{Exercise~\ref{ex:c-closure-5}: the positive-mode bootstrap.}
\phantomsection
\label{sol:ex:c-closure-5}
For a weight-two state, the Virasoro commutator formula gives
\[
 [L_1,Y(a,z)]=z^2\partial_zY(a,z)+4zY(a,z)+Y(L_1a,z).
\]
Comparing the coefficient of $z^{-n-1}$ yields
\[
 [L_1,a_{(n)}]=(2-n)a_{(n+1)}+(L_1a)_{(n)}.
\]
Quasiprimary means $L_1a=0$, so the last term vanishes. Apply
the $n=0$ identity to $p$ and rearrange:
\[
 2a_{(1)}p=L_1(a_{(0)}p)-a_{(0)}L_1p.
\]
Thus control of the zero product, together with the already proved
$L_1$ stability of its output space, controls the next positive
mode. Quasiprimary does not mean that $L_2a$ vanishes; this is why
the subsequent $L_2$ argument separately keeps track of scalar terms.



\paragraph{Exercise~\ref{ex:c-closure-6}: why one induction index fails.}
\phantomsection
\label{sol:ex:c-closure-6}
In the term $b_{(n)}(a_{(0)}t)$, a normal expression for
$a_{(0)}t$ may contain a single generator of weight $\wt(t)+1$.
The left input $b$ has weight $\wt(s)-1$, so the total input
weight is unchanged. Induction on total weight alone has made no
progress. The left weight has decreased, however, so lexicographic
induction applies. Composite and derivative terms are already dealt
with by filtered Wick and do not need this second index.


